\section{Soberanía de datos y Data Embassy: ubicación, control y recuperación son problemas distintos}

Los sistemas de alcance nacional enfrentan la data sovereignty. En la clase se mencionaron la data embassy y la virtualización, pero este tipo de conceptos se simplifica con facilidad en «los datos están dentro o fuera del país». Esta sección trata solo la dimensión arquitectónica y no deriva de la entrevista conclusiones jurídicas concretas: location, control, jurisdiction, portability y continuity deben diseñarse por separado y verificarse en conjunto mediante contratos, controles técnicos y simulacros.

\subsection{Cinco preguntas en lugar de una location checkbox}

\term{data sovereignty} (soberanía de datos) describe a qué leyes nacionales, derechos de control y políticas están sujetos los datos; \term{data residency} (residencia de datos) es más acotada y se ocupa de la ubicación física o lógica donde se almacenan los datos; \term{data embassy} (embajada de datos) suele referirse a brindar protección soberana y continuidad a activos digitales críticos mediante acuerdos jurídicos especiales e infraestructura transfronteriza. No es un producto estándar de uso general; es necesario leer los tratados y textos legales específicos.

\lecturefigure{13-data-sovereignty.png}{La soberanía de datos exige diseñar a la vez location, control, jurisdiction y continuity.}{Subtítulos locales 32:20--36:50; antecedentes del material de Saudi DGA cloud-governance.}

\begin{knowledgebox}{Lectura de la figura: las cinco preguntas de la revisión de arquitectura}
¿Dónde se almacenan en reposo y dónde se procesan los datos? ¿Quién posee las keys y la privileged identity? ¿Qué jurisdicción puede exigir una disclosure? ¿Pueden los datos y los modelos someterse a un portable export? ¿Cómo se recupera el sistema ante una falla del provider, de la region o de la red nacional? Solo cuando las cinco preguntas tienen un owner y un test claramente definidos, la sovereignty deja de ser un eslogan y se convierte en un sistema de control.
\end{knowledgebox}

\begin{warningbox}{La ubicación no aporta por sí sola seguridad ni disponibilidad}
La localización en una sola región puede aumentar las correlated failures; la replicación en varias regiones, a su vez, puede introducir problemas transfronterizos y de consistencia. Encryption, key custody, identity, logging, supply-chain assurance, backup integrity y restore drill deben diseñarse junto con la location policy.
\end{warningbox}

\teachervoice{El valor de la clase está en tratar la resiliencia nacional como un infrastructure requirement y no como un cumplimiento normativo a posteriori. Estas notas agregan lo siguiente: cualquier propuesta de «embajada digital» debe traducir los compromisos jurídicos en key ownership, replication topology, failover authority y simulacros periódicos.}

\subsection{Resumen de la sección}

La data sovereignty es un problema de arquitectura multidimensional. La residency es solo uno de sus insumos; el control, la jurisdicción, la portability, la recuperación y la auditoría determinan si el sistema está realmente disponible en una crisis.
